%!TEX root = ../thesis.tex
% Thread and channel layout for \cref{sec:impl-efficientzero}: the four worker
% kinds, the one piece of shared mutable state they contend for, and the learner
% that drives them phase by phase. The algorithm-level counterpart is
% \cref{fig:ez-workers}, which draws the same loop as data flow; this one draws
% what is actually spawned and how the threads talk to each other.
%
% Layout rules:
%  * Three bands. The middle band is the pipeline in execution order, left to
%    right --- data workers, replay buffer, rollout workers, reanalyze workers
%    --- so every stage-to-stage edge is a short horizontal arrow between
%    neighbours. The GPU workers are the band above (every device request goes
%    straight up into it), the learner the band below (every task and status goes
%    straight down).
%  * Putting the replay buffer *inside* the row rather than into a band of its
%    own is what makes the figure crossing-free: it leaves exactly one
%    non-neighbour edge, the priority write-back from the reanalyze workers,
%    which is routed through the top lane at y=1.4 --- and the one column it
%    flies over, the rollout workers, is the one column with no GPU edge. Move
%    the buffer out of the row and that edge has nowhere left to go.
%  * Columns are 3.8cm apart and boxes are 2.6cm wide, so an edge label sitting
%    in a gap has 1.2cm (34pt) of room: at \scriptsize that is about nine
%    characters per line, which is why the in-row labels are stacked and terse.
%    "updated priorities" is single-line because it rides the long detour lane
%    rather than a gap. The picture is 14cm wide against a 14.56cm \textwidth
%    --- widen the columns before lengthening any in-row label.
%  * Blue = a worker kind, and there are deliberately exactly four of them.
%    Grey = shared state, white = the orchestrating thread.
%  * Solid arrows carry data or requests, dashed arrows carry parameters.
%  * The phase divider sits at x=5.7, which is exactly where the "sampled
%    positions" label already sits, and that label is wider than the 34pt gap
%    it nominally occupies. So the divider is two segments with the label in
%    between rather than one line behind a fill=white label: a fill wide enough
%    to hide the line also reaches into the two boxes either side, and one
%    anchored tight enough not to would eat the arrow it labels. The upper
%    segment starts at samplbl.north so it tracks the label if it is reworded.
%  * The divider is drawn after the pipeline arrows but before the priority
%    lane, so the lane crosses over it rather than under.
\begin{figure}[htbp]
  \centering
  \begin{tikzpicture}[
      >={Stealth[round]}, font=\small,
      worker/.style = {draw, rounded corners=2pt, fill=blue!8,
                       minimum height=15mm, minimum width=26mm, align=center},
      store/.style  = {draw, fill=black!8,
                       minimum height=15mm, minimum width=26mm, align=center},
      band/.style   = {draw, rounded corners=2pt, fill=blue!8, align=center},
      engine/.style = {draw, rounded corners=1.5pt, fill=white, align=center,
                       font=\scriptsize, minimum width=56mm, minimum height=11mm},
      orch/.style   = {draw, rounded corners=2pt, fill=white, align=center,
                       minimum width=140mm, minimum height=12mm},
      flow/.style   = {->, thick, black!70},
      param/.style  = {->, thick, black!70, densely dashed},
      lbl/.style    = {font=\scriptsize, inner sep=2.5pt, align=center},
    ]
    % --- the device band: the only threads that touch the GPU ----------------
    \node[band, minimum width=140mm, minimum height=20mm] (gpu) at (5.7,3.6) {};
    \node at (5.7,4.28) {GPU workers};
    \node[engine] (inf)   at (2.4,3.25) {model inference\\
                                         read-only model snapshot};
    \node[engine] (train) at (9.0,3.25) {model training\\
                                         autodiff, optimizer, gradient steps};
    \draw[param] (train.west) -- node[lbl,above] {$\theta$} (inf.east);

    % --- the pipeline, in execution order ------------------------------------
    \node[worker] (data) at (0,0)    {Data\\workers\\[1pt]
                                      \scriptsize self-play episodes};
    \node[store]  (buf)  at (3.8,0)  {Replay\\buffer\\[1pt]
                                      \scriptsize Mutex-guarded};
    \node[worker] (roll) at (7.6,0)  {Rollout\\workers\\[1pt]
                                      \scriptsize sample batches};
    \node[worker] (rean) at (11.4,0) {Reanalyze\\workers\\[1pt]
                                      \scriptsize refresh targets};

    \draw[flow] (data) -- node[lbl,above] {finished\\episodes}  (buf);
    \draw[flow] (buf)  -- node[lbl,above] (samplbl) {sampled\\positions} (roll);
    \draw[flow] (roll) -- node[lbl,above] {raw\\batches}        (rean);

    % --- the phase divider ---------------------------------------------------
    % Cuts the row between the replay buffer and the rollout workers, which is
    % where the two phases actually part: everything from the rollout workers
    % rightwards runs only in the training phase. It stops short of the GPU band
    % and the learner --- both are live in both phases, so extending it through
    % them would claim a split that does not exist.
    \draw[black!35, densely dashed] (5.7,-1.6) -- (5.7,0);
    \draw[black!35, densely dashed] (samplbl.north) -- (5.7,2.15);
    \node[lbl, black!55, anchor=south east] at (5.6,2.2) {data phase};
    \node[lbl, black!55, anchor=south west] at (5.8,2.2) {training phase};

    % --- device requests: straight up, one per column that needs one ---------
    \draw[flow] (data.north) --
      node[lbl,right,pos=0.62] {batched\\MCTS} (data.north |- gpu.south);
    \draw[flow] ([xshift=7mm]rean.north) --
      node[lbl,left,pos=0.7] {target \gls{mcts},\\gradient step}
      ([xshift=7mm]rean.north |- gpu.south);

    % --- the one non-neighbour edge, routed through the top lane -------------
    \draw[flow,rounded corners=3pt]
      ([xshift=-7mm]rean.north) -- (10.7,1.4)
      -- node[lbl,fill=white] {updated priorities} (3.8,1.4) -- (buf.north);

    % --- the orchestrator: tasks down, statuses up ---------------------------
    \node[orch] (learner) at (5.7,-2.35) {Learner\\[1pt]
      \scriptsize one iteration = a data phase, then a training phase};

    \draw[<->,thick,black!70] (data.south) --
      node[lbl,right] {tasks,\\statuses} (learner.north -| data.south);
    \draw[flow] (learner.north -| roll.south) --
      node[lbl,right] {tasks} (roll.south);
    \draw[flow] (rean.south) --
      node[lbl,right] {statuses} (learner.north -| rean.south);
  \end{tikzpicture}
  \caption[The worker threads of the EfficientZero implementation]{%
    The workers spawned by the learner and how they interact.
    Data workers receive data generation tasks and save episodes in the replay buffer.
    Rollout workers receive rollout tasks, sample raw batches from the replay buffer
    and forward them to reanalyze workers.
    Reanalyze workers refresh the raw batches using the latest model and then use
    backpropagation to calculate losses and step the model optimizer.
    The dashed line separates the data collection and training phases of a single iteration.
  }
  \label{fig:impl-ez-workers}
\end{figure}
